\documentclass[12pt]{article}
\usepackage[utf8]{inputenc}
\usepackage[T1]{fontenc}
\usepackage{hyperref}
\title{Parsi Language Resources for K-12 Curriculum Development}
\author{Compiled by Claude Sonnet 5.5 (claude-sonnet-5-5) for Andr\'as Kornai}
\date{}
\begin{document}
\maketitle

\section{Introduction}
The code \texttt{prp} (Parsi) was deprecated and merged into Gujarati (\texttt{guj}) in ISO 639-3, effective 2023-01-20. Glottolog (\texttt{pars1251}) retains the entry only for bookkeeping and treats Parsi as the variety of Gujarati spoken by the Parsis, not a distinct language. It is also distinct from Parsi-Dari (\texttt{prd}), a supposed Zoroastrian variety in Iran, which Glottolog treats as spurious. Resources for the Parsi variety therefore fall under Gujarati. I did not search for Gujarati school materials, so this dossier lists only what bears on the identification.

\section{Primary Educational Resources}
None searched. See Gujarati.

\section{Linguistic and Reference Materials}
\begin{itemize}
\item Glottolog entry (retired): \url{https://glottolog.org/resource/languoid/id/pars1251}
\item Wikipedia, Parsi language: \url{https://en.wikipedia.org/wiki/Parsi_language}
\item Wikipedia, Gujarati languages: \url{https://en.wikipedia.org/wiki/Gujarati_languages}
\end{itemize}

\section{Religious and Community Resources}
None found.

\section{Digital and Audio Resources}
None found.

\section{Curriculum Development Context}
A separate Parsi curriculum is not supported by the sources found. The practical route is the Gujarati curriculum, with community-specific vocabulary and heritage content added if wanted.

\end{document}